% Part: normal-modal-logic
% Chapter: sequent-calculus
% Section: rules-for-K

\documentclass[../../../include/open-logic-section]{subfiles}

\begin{document}

\olfileid{nml}{seq}{rul}

\olsection{Aturan kanggo \Log{K}}

Aturan kanggo konektif proposisional lumrah padha karo aturan ing
kalkulus sekuen lumrah~$\Log{LK}$. Aksiomane uga padha: saben sekuen
awujud $!A \Sequent !A$ dianggep aksioma.

Kanggo operator modal\iftag{prvBox}{\iftag{prvDiamond}{~$\Box$
lan}{~$\Box$}}{}\iftag{prvDiamond}{~$\Diamond$}{}, ana
\iftag{notprvBox,notprvDiamond}{aturan}{aturan} tambahan iki:
\[
  \iftag{prvBox}{\iftag{prvDiamond}{% <> and [] primitive
    \Axiom$\Gamma \fCenter \Delta, !A$
    \RightLabel{$\Box$}
    \UnaryInf$\Box\Gamma \fCenter \Diamond\Delta, \Box!A$
    \DisplayProof  
    \qquad
    \Axiom$!A, \Gamma \fCenter \Delta$
    \RightLabel{$\Diamond$}
    \UnaryInf$\Diamond!A, \Box\Gamma \fCenter \Diamond\Delta$
    \DisplayProof
}{% only Box primitive
\Axiom$\Gamma \fCenter !A$
\RightLabel{$\Box$}
\UnaryInf$\Box\Gamma \fCenter \Box!A$
\DisplayProof
}}{% only <> primitive
  \Axiom$!A \fCenter \Delta$
  \RightLabel{$\Diamond$}
  \UnaryInf$\Diamond!A \fCenter \Diamond\Delta$
  \DisplayProof
}
\]
Ing kene, \iftag{prvBox}{$\Box\Gamma$ ateges urutan !!{formula}s
kang diasilake saka~$\Gamma$ kanthi masang $\Box$ ing ngarepe saben
!!{formula} ing~$\Gamma$\iftag{prvDiamond}{ lan
}{}}{}\iftag{prvDiamond}{$\Diamond\Delta$ ateges urutan
!!{formula}s kang diasilake saka~$\Delta$ kanthi masang $\Diamond$
ing ngarepe saben !!{formula} ing~$\Delta$}{}.
\iftag{notprvDiamond}{Sisih tengen premis aturan $\Box$ mung kena
ngemot paling akeh siji !!{formula}~$!A$. }{}%
\iftag{notprvBox}{Sisih kiwa premis aturan $\Diamond$ mung kena
ngemot paling akeh siji !!{formula}~$!A$. }{}%
\iftag{prvBox}{$\Gamma$\iftag{prvDiamond}{~lan~}{}}{}\iftag{prvDiamond}{$\Delta$}{}
bisa kosong; yen mangkono, bagean kang cocog, yaiku
\iftag{prvBox}{$\Box\Gamma$\iftag{prvDiamond}{~lan~}{}}{}\iftag{prvDiamond}{$\Diamond\Delta$}{}
ing sekuen kesimpulan uga kosong.

Watesan yen
\iftag{prvBox}{$\Box$ ing sisih tengen\iftag{prvDiamond}{ lan
}{}}{}\iftag{prvDiamond}{$\Diamond$ ing sisih kiwa}{}
mung ditambahake marang siji !!{formula}~$!A$ iku perlu. Yen kita
ngolehake
\iftag{prvBox}{nambahake $\Box$ marang !!{formula}s saakeh-akehe
ing sisih tengen\iftag{prvDiamond}{ utawa }{}}{}\iftag{prvDiamond}{nambahake $\Diamond$
marang !!{formula}s saakeh-akehe ing sisih kiwa}{}, kita bisa
!!{derive}:
  \[
    \iftag{prvBox}{
      \Axiom$!A \fCenter !A$
      \RightLabel{\RightR{\lnot}}
      \UnaryInf$\fCenter !A, \lnot !A$
      \RightLabel{$\Box*$}
      \UnaryInf$\fCenter \Box!A, \Box\lnot !A$
      \doubleLine
      \RightLabel{\RightR{\lor}}
      \UnaryInf$\fCenter \Box!A \lor \Box \lnot !A$
      \DisplayProof}{}
    \iftag{notprvBox,notprvDiamond}{}{\qquad}
    \iftag{prvDiamond}{
      \Axiom$!A \fCenter !A$
      \RightLabel{\LeftR{\lnot}}
      \UnaryInf$\lnot !A, !A \fCenter$
      \RightLabel{$\Diamond*$}
      \UnaryInf$\Diamond\lnot !A,\Diamond!A \fCenter $
      \RightLabel{\RightR{\lnot}}
      \UnaryInf$\Diamond!A \fCenter \lnot\Diamond\lnot !A$
      \RightLabel{\RightR{\lif}}
      \UnaryInf$ \fCenter \Diamond!A \lif \lnot\Diamond\lnot !A$
      \DisplayProof}{}
  \]
Nanging, \iftag{prvBox}{$\Box!A \lor \Box \lnot !A$\iftag{prvDiamond}{ lan
$\Diamond!A \lif \lnot\Diamond\lnot !A$}{}}{$\Diamond!A \lif
\lnot\Diamond\lnot !A$} ora valid ing~$\Log{K}$.

Yen kita ngolehake rumus sampingan saliyane~$!A$ ing premis lan
ngolehake \iftag{prvBox}{aturan $\Box$ nambahake $\Box$ mung marang~$!A$
ing sisih tengen\iftag{prvDiamond}{, utawa ngolehake }{}}{}\iftag{prvDiamond}{aturan
$\Diamond$ nambahake $\Diamond$ mung marang~$!A$ ing sisih kiwa}{}
(tanpa ngowahi rumus sampingan), kita bisa !!{derive}:
\[
  \iftag{prvBox}{
    \Axiom$!A \fCenter !A$
    \RightLabel{\RightR{\lnot}}
    \UnaryInf$\fCenter !A, \lnot !A$
    \RightLabel{\RightR{\Exchange}}
    \UnaryInf$\fCenter \lnot !A, !A$
    \RightLabel{$\Box*$}
    \UnaryInf$\fCenter \lnot!A, \Box !A$
    \doubleLine
    \RightLabel{\RightR{\lor}}
    \UnaryInf$\fCenter \lnot!A \lor \Box !A$
    \DisplayProof}{}
  \iftag{notprvBox,notprvDiamond}{}{\qquad}
  \iftag{prvDiamond}{
    \Axiom$!A \fCenter !A$
    \RightLabel{\LeftR{\lnot}}
    \UnaryInf$\lnot !A, !A \fCenter$
    \RightLabel{$\Diamond*$}
    \UnaryInf$\Diamond\lnot !A, !A \fCenter $
    \RightLabel{\RightR{\lnot}}
    \UnaryInf$!A \fCenter \lnot\Diamond\lnot !A$
    \RightLabel{\RightR{\lif}}
    \UnaryInf$ \fCenter !A \lif \lnot\Diamond\lnot !A$
    \DisplayProof}{}
\]
Nanging, \iftag{prvBox}{$\lnot!A \lor \Box !A$ (kang ekuivalen karo $!A
\lif \Box!A$)\iftag{prvDiamond}{ lan $!A \lif \lnot\Diamond\lnot !A$
}{}}{$!A \lif \lnot\Diamond\lnot !A$} ora valid ing~$\Log{K}$.

\end{document}
